\section{Sets: The Grammar of Modern Mathematics}
\chapterline

In elementary mathematics, arithmetic operates on numbers, algebra operates on equations, and calculus operates on functions. But when you ask what numbers, equations, and functions actually \emph{are}, modern mathematics answers in a single unified vocabulary: they are all built from \textbf{sets}.

A set is a well-defined collection of distinct objects. The objects belonging to a set are called its \textbf{elements} (or members). We write:
\begin{align*}
x \in A &\quad \text{to assert that } x \text{ is an element of } A, \\
x \notin A &\quad \text{to assert that } x \text{ is not an element of } A.
\end{align*}

This foundation seems deceptively simple. Yet sets serve as the universal scaffolding for advanced mathematics: spaces are sets of points equipped with distance or topology; groups are sets equipped with lawful operations; relations are sets of ordered pairs; and differential equations are conditions selecting subsets of function spaces.

\begin{historybox}
In the 1870s, the German mathematician Georg Cantor invented set theory while studying which collections of numbers allow a trigonometric wave to be uniquely reconstructed. Cantor realized that to tame infinity, mathematics needed a language that could treat infinite collections as completed objects in their own right. He showed that sets could serve as the universal atom out of which numbers, functions, and geometric spaces are built.

At the turn of the twentieth century, Gottlob Frege attempted to ground all of arithmetic and mathematics in this new foundation in his monumental \emph{Grundgesetze der Arithmetik}. Frege relied on what seemed like an obvious truth: \emph{Unrestricted Comprehension}---the belief that any mathematical condition you can write down creates a valid set of all objects satisfying it. But in June 1902, the British philosopher Bertrand Russell sent Frege a brief letter that shattered naive set theory. Russell asked: what happens if you consider the set of all sets that do not contain themselves?
\[
R = \{x : x \notin x\}.
\]
If $R$ is not an element of itself, it meets the entrance criteria, so it must belong to $R$. But if $R$ does belong to $R$, it fails the condition, so it cannot belong to $R$! The foundation of logic had generated an inescapable contradiction. Frege wrote in an appendix that a scientist can hardly meet with anything more undesirable than having their foundation collapse just as the work is finished.

Mathematicians led by Ernst Zermelo (1908) and Adolf Fraenkel (1922) repaired the foundation by replacing naive set theory with modern axiomatic set theory (ZFC). Under the new rules, you can only carve out subsets from an \emph{already existing} set ($\{x \in A : P(x)\}$), banishing self-referential monsters and the dangerous ``set of all sets.'' Sets became the bedrock of modern mathematics.
\end{historybox}

\subsection{The crucial distinction: element versus subset}

The single most common source of early confusion is conflating \textbf{membership} ($\in$) with \textbf{containment} ($\subseteq$).

\begin{itemize}
\item \textbf{Membership ($x \in A$)}: A relation between an individual object and a container. $x$ is one of the items sitting inside the box $A$.
\item \textbf{Subset ($A \subseteq B$)}: A relation between two containers. $A \subseteq B$ means that \emph{every} element of $A$ is also an element of $B$:
\[
A \subseteq B \iff \forall x, \, (x \in A \implies x \in B).
\]
\end{itemize}

To see how easy it is to trip if you are not careful, consider the set:
\[
S = \bigl\{1, \, 2, \, \{3, 4\}\bigr\}.
\]
Notice that $S$ contains exactly \emph{three} elements: the number $1$, the number $2$, and the package $\{3, 4\}$.
\begin{itemize}
\item $1 \in S$ is \textbf{True}.
\item $\{1\} \subseteq S$ is \textbf{True} (because $1 \in S$).
\item $\{1\} \in S$ is \textbf{False} (the number $1$ is inside $S$, but the set $\{1\}$ is not).
\item $\{3, 4\} \in S$ is \textbf{True} (the package $\{3, 4\}$ is literally an element of $S$).
\item $\{3, 4\} \subseteq S$ is \textbf{False} (because $3 \notin S$ and $4 \notin S$; the elements $3$ and $4$ are trapped inside the inner box).
\item $\bigl\{\{3, 4\}\bigr\} \subseteq S$ is \textbf{True} (a set containing the element $\{3, 4\}$).
\end{itemize}

Think of a box containing an apple, an orange, and a sealed envelope with two photographs inside. The apple is an element of the box. The envelope is an element of the box. But the photographs are \emph{not} loose elements of the box; they are elements of the envelope!

\subsection{The four fundamental set operations}

Given two sets $A$ and $B$ living inside a universal set $U$, we can combine and filter them using four primary operations:

\begin{center}
\begin{tikzpicture}[
    font=\small,
    subpanel/.style={draw=gray!30, fill=gray!2, rounded corners=4pt, inner sep=6pt}
]

% Panel 1: Intersection (Top Left)
\node[subpanel, text width=5.8cm, minimum height=3.2cm] (p1) at (0, 3.8) {};
\node[anchor=west, font=\footnotesize\scshape, text=sourceblue!90!black] at (-2.7, 5.1) {Intersection: $A \cap B$};
\begin{scope}[xshift=-0.5cm, yshift=3.6cm, scale=0.75]
    \draw[thick, draw=sourceblue!70, fill=white] (-0.8, 0) circle (1.2);
    \draw[thick, draw=targetgreen!70, fill=white] (0.8, 0) circle (1.2);
    \begin{scope}
        \clip (-0.8, 0) circle (1.2);
        \fill[sourceblue!45] (0.8, 0) circle (1.2);
    \end{scope}
    \node[font=\scriptsize] at (-1.3, 0) {$A$};
    \node[font=\scriptsize] at (1.3, 0) {$B$};
\end{scope}
\node[font=\scriptsize, text=gray!80!black] at (0, 2.5) {$\{x : x \in A \land x \in B\}$ (Elements in both)};

% Panel 2: Union (Top Right)
\node[subpanel, text width=5.8cm, minimum height=3.2cm] (p2) at (6.5, 3.8) {};
\node[anchor=west, font=\footnotesize\scshape, text=targetgreen!90!black] at (3.8, 5.1) {Union: $A \cup B$};
\begin{scope}[xshift=6.0cm, yshift=3.6cm, scale=0.75]
    \draw[thick, draw=sourceblue!70, fill=targetgreen!30] (-0.8, 0) circle (1.2);
    \draw[thick, draw=targetgreen!70, fill=targetgreen!30] (0.8, 0) circle (1.2);
    \node[font=\scriptsize] at (-1.3, 0) {$A$};
    \node[font=\scriptsize] at (1.3, 0) {$B$};
\end{scope}
\node[font=\scriptsize, text=gray!80!black] at (6.5, 2.5) {$\{x : x \in A \lor x \in B\}$ (Inclusive OR)};

% Panel 3: Set Difference (Bottom Left)
\node[subpanel, text width=5.8cm, minimum height=3.2cm] (p3) at (0, 0) {};
\node[anchor=west, font=\footnotesize\scshape, text=quotientorange!90!black] at (-2.7, 1.3) {Set Difference: $A \setminus B$};
\begin{scope}[xshift=-0.5cm, yshift=-0.2cm, scale=0.75]
    \draw[thick, draw=sourceblue!70, fill=quotientorange!40] (-0.8, 0) circle (1.2);
    \draw[thick, draw=targetgreen!70, fill=white] (0.8, 0) circle (1.2);
    \begin{scope}
        \clip (0.8, 0) circle (1.2);
        \fill[white] (-0.8, 0) circle (1.2);
    \end{scope}
    \draw[thick, draw=sourceblue!70] (-0.8, 0) circle (1.2);
    \draw[thick, draw=targetgreen!70] (0.8, 0) circle (1.2);
    \node[font=\scriptsize] at (-1.3, 0) {$A$};
    \node[font=\scriptsize] at (1.3, 0) {$B$};
\end{scope}
\node[font=\scriptsize, text=gray!80!black] at (0, -1.3) {$\{x : x \in A \land x \notin B\}$ (In $A$, stripped of $B$)};

% Panel 4: Symmetric Difference (Bottom Right)
\node[subpanel, text width=5.8cm, minimum height=3.2cm] (p4) at (6.5, 0) {};
\node[anchor=west, font=\footnotesize\scshape, text=structurepurple!90!black] at (3.8, 1.3) {Symmetric Difference: $A \triangle B$};
\begin{scope}[xshift=6.0cm, yshift=-0.2cm, scale=0.75]
    \draw[thick, draw=sourceblue!70, fill=structurepurple!35] (-0.8, 0) circle (1.2);
    \draw[thick, draw=targetgreen!70, fill=structurepurple!35] (0.8, 0) circle (1.2);
    \begin{scope}
        \clip (-0.8, 0) circle (1.2);
        \fill[white] (0.8, 0) circle (1.2);
    \end{scope}
    \draw[thick, draw=sourceblue!70] (-0.8, 0) circle (1.2);
    \draw[thick, draw=targetgreen!70] (0.8, 0) circle (1.2);
    \node[font=\scriptsize] at (-1.3, 0) {$A$};
    \node[font=\scriptsize] at (1.3, 0) {$B$};
\end{scope}
\node[font=\scriptsize, text=gray!80!black] at (6.5, -1.3) {$(A \setminus B) \cup (B \setminus A)$ (Exclusive OR / XOR)};

\end{tikzpicture}
\end{center}

\subsection{Cartesian products: building multidimensional worlds}

When we want to pair elements together in a specific sequence, set membership alone is insufficient because sets are unordered ($\{1, 2\} = \{2, 1\}$). To introduce order, mathematics defines the \textbf{ordered pair} $(a, b)$, where order matters: $(1, 2) \neq (2, 1)$.

\begin{definition}
The \emph{Cartesian product} of two sets $A$ and $B$, denoted $A \times B$, is the set of all ordered pairs whose first coordinate comes from $A$ and second coordinate comes from $B$:
\[
A \times B = \{(a, b) : a \in A \text{ and } b \in B\}.
\]
\end{definition}

If $A = \{1, 2, 3, 4\}$ and $B = \{u, v, w\}$, the Cartesian product $A \times B$ forms a discrete $4 \times 3$ coordinate grid:

\begin{center}
\begin{tikzpicture}[
    font=\small,
    axis/.style={draw=gray!70, -{Latex[length=2.5mm]}, thick},
    gridpt/.style={circle, fill=sourceblue, inner sep=2.2pt}
]

% Background grid panel
\draw[draw=gray!25, fill=gray!2, rounded corners=4pt] (-0.8, -0.6) rectangle (8.0, 4.3);
\node[anchor=north west, font=\footnotesize\scshape, text=sourceblue!90!black] at (-0.6, 4.15) {Coordinate Grid of $A \times B$};

% Coordinate axes
\draw[axis] (0.5, 0.4) -- (7.2, 0.4) node[right=4pt, font=\scriptsize] {Set $A$};
\draw[axis] (0.5, 0.4) -- (0.5, 3.2) node[above=4pt, font=\scriptsize] {Set $B$};

% Ticks on A
\foreach \x/\label in {2.0/1, 3.5/2, 5.0/3, 6.5/4} {
    \draw[thick, gray!70] (\x, 0.3) -- (\x, 0.5) node[below=4pt, font=\scriptsize\bfseries] {$\label$};
}

% Ticks on B
\foreach \y/\label in {1.2/u, 2.0/v, 2.8/w} {
    \draw[thick, gray!70] (0.4, \y) -- (0.6, \y) node[left=4pt, font=\scriptsize\bfseries] {$\label$};
}

% Highlighted slice: {2} x B (drawn behind grid points so labels stay crisp)
\draw[draw=targetgreen!80, fill=targetgreen!12, rounded corners=3pt, thick] (3.2, 0.9) rectangle (4.35, 3.25);
\node[font=\tiny\bfseries, text=targetgreen!90!black, align=center] at (3.77, 3.5) {Fiber: $\{2\} \times B$};

% Grid points
\foreach \x/\lx in {2.0/1, 3.5/2, 5.0/3, 6.5/4} {
    \foreach \y/\ly in {1.2/u, 2.0/v, 2.8/w} {
        \draw[dotted, gray!40] (\x, 0.4) -- (\x, 2.8);
        \draw[dotted, gray!40] (0.5, \y) -- (6.5, \y);
        \node[gridpt] at (\x, \y) {};
        \node[above right=1pt, font=\tiny, text=gray!80!black] at (\x, \y) {$(\lx, \ly)$};
    }
}

\end{tikzpicture}
\end{center}

When both sets are the real line $\R$, the Cartesian product $\R \times \R = \R^2$ is the familiar Euclidean plane. Every point in coordinate geometry is simply an element of a Cartesian product!

\subsection{Proving set equality: double containment}

How do you prove that two set expressions describe the exact same set, such as $A \cap (B \cup C) = (A \cap B) \cup (A \cap C)$?

You cannot merely draw a Venn diagram and declare it finished. A formal proof of $S = T$ requires establishing \textbf{mutual containment} in two separate phases:
\begin{enumerate}
\item \textbf{Show $S \subseteq T$}: Let $x \in S$ be arbitrary. Trace $x$ through the definitions until you deduce $x \in T$.
\item \textbf{Show $T \subseteq S$}: Let $x \in T$ be arbitrary. Trace $x$ through the definitions until you deduce $x \in S$.
\end{enumerate}

\begin{theorem}
For any sets $A$ and $B$, $A \cap B = A \setminus (A \setminus B)$.
\end{theorem}

\begin{proof}
We prove both inclusions:

\textbf{Direction 1: Show $A \cap B \subseteq A \setminus (A \setminus B)$.}
Let $x \in A \cap B$ be arbitrary. By definition of intersection, $x \in A$ and $x \in B$.
Now examine the set $A \setminus B = \{y : y \in A \land y \notin B\}$. Since $x \in B$, $x$ cannot satisfy the condition $x \notin B$. Thus, $x \notin (A \setminus B)$.
Since we know $x \in A$ and $x \notin (A \setminus B)$, by definition of set difference we have $x \in A \setminus (A \setminus B)$. This establishes the first inclusion.

\textbf{Direction 2: Show $A \setminus (A \setminus B) \subseteq A \cap B$.}
Let $x \in A \setminus (A \setminus B)$ be arbitrary. By definition of set difference, $x \in A$ and $x \notin (A \setminus B)$.
Recall that $A \setminus B$ is the set of elements in $A$ that are not in $B$. Since $x \in A$, the only way for $x$ to fail to belong to $A \setminus B$ is if $x \in B$.
Since $x \in A$ and $x \in B$, by definition of intersection we have $x \in A \cap B$. This establishes the second inclusion.

Having established both containments, we conclude that $A \cap B = A \setminus (A \setminus B)$.
\end{proof}

\subsection{What set grammar unlocked---and the need for dynamic maps}

Set theory unlocked a universal language for the mathematical sciences. Concepts from disparate fields---vectors in physics, events in probability, algorithms in computer science, and geometric surfaces in calculus---are all formalized in the single unified language of sets, membership, subsets, and Cartesian coordinate grids. Proving mutual containment through element chasing provides an unshakeable algorithmic technique for establishing that two complex mathematical structures are identical.

Yet sets alone have an honest limitation: \textbf{they are purely static collections}. A set tells you which elements live inside its borders, but it does not describe motion, transformation, or structural relationships between different worlds. How do you map a point from the Euclidean plane into a curved sphere? How do you transform a vector without tearing space? 

To breathe dynamic life into static collections, we need structured mappings. In Chapter~11 (\emph{Set Proofs: Element Chasing in Detail}), we master the fine art of element-chasing proofs, and in Chapter~12 (\emph{Functions: Structured Promises, Not Just Formulas}), we build the modern theory of functions that connect sets across the mathematical universe.

\subsection{Exercises}

Let $A = \{1, 2, 3\}$, $B = \{3, 4\}$, and $C = \{1, 4\}$.
\begin{enumerate}
\item Compute $A \cup B$, $A \cap B$, $A \setminus B$, and $A \triangle B$.
\item Write down all elements of the Cartesian product $B \times C$.
\item Write down the power set $\pow(B)$. How many elements does it have?
\item True or false? Explain your reasoning:
\begin{enumerate}
\item $\emptyset \in \pow(A)$.
\item $\emptyset \subseteq \pow(A)$.
\item $\{\emptyset\} \subseteq \pow(A)$.
\item $\{1, 2\} \in \pow(A)$.
\item $\{1, 2\} \subseteq \pow(A)$.
\end{enumerate}
\item Prove via element chasing that for any sets $A, B, C$:
\[
A \setminus (B \cup C) = (A \setminus B) \cap (A \setminus C).
\]
\end{enumerate}

\newpage
